% Source-faithful assembly: Mumford, Abelian Varieties, Chapter III,
% Horizon transcription pages 0132--0176 (with quotient continuation pages
% 0127--0131 assembled in ch03-source-part1.tex).

\section{Finite-fibre associated spaces}
\label{mumford-III12-associated-source}
\dcref{M:ch3:III.12}

\begin{definition}[Fibration with finite fibre]
\label{mumford-def-associated-finite-fibre}
\dcref{M:ch3:III.12}
Let a finite group scheme $G$ act freely on $X$, with quotient $X/G$, and
let $F$ be a finite scheme with a $G$-action. The induced free action on
$X\times F$ has quotient
\[
U=(X\times F)/G,
\]
and the natural morphism $\pi:U\to V=X/G$ is called the fibration with
fibre $F$ associated to the principal $G$-bundle $X\to X/G$. It is finite
and flat; its degree is the rank of $\pi_*\mathcal O_U$ over
$\mathcal O_V$.
\source{mumford-abelian-varieties:page-0132}
\uses{mumford-thm-finite-groupscheme-quotient}
\end{definition}

\begin{theorem}[Euler characteristic of an associated fibration]
\label{mumford-thm-euler-associated-fibration}
\dcref{M:ch3:III.12}
\uses{mumford-def-associated-finite-fibre}
If $\pi:U\to V$ is a finite-fibre fibration associated to a principal
$G$-bundle, then for every coherent sheaf $\mathcal F$ on $V$,
\[
\chi(\pi^*\mathcal F)=(\deg\pi)\chi(\mathcal F).
\]
\source{mumford-abelian-varieties:page-0132}
\source{mumford-abelian-varieties:page-0133}
\end{theorem}
\begin{proof}
It is enough to treat the case $V=U/G$ for a free action, since
$(X\times F)\to(X\times F)/G\to X/G$ reduces the general case to the
composite and first morphism. Suppose the assertion fails, and choose a
counterexample $\mathcal F$ minimal by support. Let $\mathcal J$ annihilate
$\mathcal F$ and write $V'$ for its support, $U'=\pi^{-1}(V')$. The quotient
proposition shows that $U'\to V'$ remains a principal $G$-space. Thus we
may assume the assertion holds for every coherent sheaf with proper support.
Exact sequences and flatness then show it also holds for extensions whenever
it holds for two terms.

If $V$ is reducible, a short exact sequence with both terms supported on
proper closed subsets gives a contradiction. If the nilpotent ideal
$\mathcal J$ of $V$ is nonzero, the exact sequence
\[
0\to\mathcal J\mathcal F\to\mathcal F\to
\mathcal F/\mathcal J\mathcal F\to0
\]
does the same. Hence assume $V$ reduced and irreducible. If $r$ is the
generic rank of $\mathcal F$, there is an ideal $\mathcal J$ and an
injection $\mathcal F^r\to\mathcal F$ with proper-support cokernel.
Consequently it suffices to prove the formula for $\mathcal O_V$.

Apply it to $\pi_*\mathcal O_U$, a coherent sheaf of nonzero rank. By the
projection formula,
\[
\begin{aligned}
\chi(\pi^*\pi_*\mathcal O_U)
 &=\chi(\mathcal O_V\otimes_{\mathcal O_V}\pi_*\mathcal O_U)\\
 &=(\deg\pi)\chi(\mathcal O_V)
 =(\deg\pi)\chi(\pi_*\mathcal O_U),
\end{aligned}
\]
which proves the assertion.
\end{proof}

\begin{remark}[Finite étale maps as associated spaces]
\label{mumford-rem-symmetric-associated-space}
\dcref{M:ch3:III.12}
Every finite étale morphism $f:U\to V$ of constant degree $n$ is an
associated finite-fibre fibration for a principal $\Sigma(n)$-space. In the
$n$-fold fibre product $U\times_V\cdots\times_VU$, let $P$ be the open and
closed locus of tuples with pairwise distinct entries. Then $P\to V$ is a
principal $\Sigma(n)$-space. If $F$ is the reduced $n$-point scheme with
the natural permutation action, the invariant map
\[
P\times F\longrightarrow V,
\qquad ((x_1,\ldots,x_n),i)\longmapsto x_i
\]
has quotient $V$ and identifies the associated space with $U$.
\source{mumford-abelian-varieties:page-0133}
\source{mumford-abelian-varieties:page-0134}
\uses{mumford-thm-euler-associated-fibration}
\end{remark}

\section{The dual abelian variety in any characteristic}
\label{mumford-III13-source}
\dcref{M:ch3:III.13}

\begin{definition}[Scheme-theoretic stabilizer of a line bundle]
\label{mumford-def-KL}
\dcref{M:ch3:III.13}
For a line bundle $L$ on an abelian variety $X$, put
\[
M=m^*L\otimes p_1^*L^{-1}\otimes p_2^*L^{-1}
\]
on $X\times X$. The group scheme $K(L)\subset X$ is the maximal closed
subscheme for which $M|_{K(L)\times X}$ is trivial. Equivalently, for every
scheme $S$ and $f:S\to X$, one has $f\in K(L)(S)$ exactly when
\[
T_f^*L_S\simeq L_S\otimes p_2^*N
\]
for a line bundle $N$ on $S$.
\source{mumford-abelian-varieties:page-0134}
\source{mumford-abelian-varieties:page-0135}
\end{definition}
\begin{proof}
The composite $X\times S\xrightarrow{T_f}X\times S\to X$ is
$m\circ(1_X\times f)$, so
$T_f^*L_S\simeq(1_X\times f)^*m^*L$. Restricting a putative isomorphism
to $\{0\}\times S$ identifies $N$ with $f^*L$. Hence the displayed
condition is equivalent to triviality of $(1_X\times f)^*M$, which is the
defining universal property of $K(L)$. In particular $K(L)$ is a subgroup
scheme of $X$.
\end{proof}

\begin{construction}[Dual quotient and normalized Poincare bundle]
\label{mumford-thm-dual-any-characteristic}
\label{mumford-prop-abelian-quotient}
\dcref{M:ch3:III.13}
\uses{mumford-def-KL,mumford-prop-maximal-triviality-scheme,
 mumford-thm-finite-groupscheme-quotient}
Choose an ample line bundle $L$ on $X$. Then $K(L)$ is finite. The quotient
\[
\pi:X\longrightarrow\widehat X=X/K(L)
\]
is an abelian variety. The bundle $M$ above admits a canonical
$K(L)$-linearization and descends uniquely to a normalized line bundle $P$ on
$\widehat X\times X$. Its restrictions to $\{0\}\times X$ and
$\widehat X\times\{0\}$ are trivial, and
\[
\widehat X\xrightarrow{\sim}\operatorname{Pic}^0X
\]
identifies the fibre of $P$ over $\alpha$ with the corresponding algebraically
trivial line bundle on $X$.
\source{mumford-abelian-varieties:page-0135}
\source{mumford-abelian-varieties:page-0136}
\end{construction}
\begin{proof}
For $x\in K(L)(S)$, write $T_x^*L_S\simeq L_S\otimes L_0$ with $L_0$ from
$S$. Then translation by $(x,0)$ preserves $M_S$. Any two lifts to $M_S$
differ by a global unit; the restriction map from $X_S\times_SX_S$ to
$X_S\times_S\{0\}$ is an isomorphism on units. Rigidifying along the second
factor therefore selects a unique lift, and these lifts form a group action.
Faithfully flat descent for the finite quotient gives $P$ and its uniqueness.
The map $X\to\operatorname{Pic}^0X$, $x\mapsto T_x^*L\otimes L^{-1}$, is
constant on $K(L)$ and has kernel $K(L)$, so the induced map from $\widehat X$
is the displayed isomorphism.
\end{proof}

\begin{theorem}[Universal Poincare mapping property]
\label{mumford-thm-poincare-universal-mapping}
\dcref{M:ch3:III.13}
\uses{mumford-prop-maximal-triviality-scheme,mumford-cor-structure-sheaf-cohomology}
Let $S$ be a scheme and $L$ a line bundle on $S\times X$, trivial on
$S\times\{0\}$, whose restriction to every fibre $\{s\}\times X$ lies in
$\operatorname{Pic}^0X$. There is a unique morphism $\phi:S\to\widehat X$
such that $L\simeq(\phi\times1_X)^*P$.
\source{mumford-abelian-varieties:page-0136}
\source{mumford-abelian-varieties:page-0137}
\source{mumford-abelian-varieties:page-0138}
\source{mumford-abelian-varieties:page-0139}
\source{mumford-abelian-varieties:page-0140}
\end{theorem}
\begin{proof}
Put $M=p_{23}^*P\otimes p_{13}^*L^{-1}$ on
$S\times\widehat X\times X$ and let $\Gamma_S$ be its maximal triviality
subscheme. Uniqueness and existence reduce to proving
$\Gamma_S\to S$ is an isomorphism. By localization and base change assume
$S=\operatorname{Spec}B$ is finite local and twist $L$ so its closed fibre
is trivial. The direct images along $S\times\widehat X\times X\to S\times X$
have discrete support, so Leray gives free $B$-modules
\[
H^p(S\times\widehat X\times X,M)
\simeq H^0(S\times X,R^pp_{13,*}M).
\]
The direct images along $S\times\widehat X\times X\to S\times\widehat X$
are supported at $(s,0)$ and are computed by a finite free complex
$0\to K_0\to\cdots\to K_g\to0$ over
$A=B\otimes_k\mathcal O_{0,\widehat X}$.

We use the following regular-local-ring lemma: if $\mathcal O$ is regular of
dimension $g$ and a length-$g$ complex of finite free modules has artinian
cohomology, then $H^i=0$ for $i<g$. For $g>0$, choose
$x\in\mathfrak M\setminus\mathfrak M^2$, reduce the complex modulo $x$,
and use the long exact cohomology sequence and induction on $g$; Nakayama
then gives the vanishing. Applying the lemma and its dual to the complex
above gives an artinian terminal module $K\simeq A/\mathfrak Q$, with
$\mathfrak Q\cap B=0$. For every $\mathfrak M$-primary ideal $\mathfrak b$,
\[
H^0(V(\mathfrak b)\times X,M)
\simeq\operatorname{Hom}_A(A/\mathfrak Q,A/\mathfrak b).
\]
The evaluation map is surjective after reduction at the closed point, hence
$M$ is trivial over $V(\mathfrak Q)\times X$ and
$\Gamma_S=V(\mathfrak Q)$. Its fibre is the reduced point $(s,0)$, so
$A/(\mathfrak Q+\mathfrak M_A)=k$ and $B\to A/\mathfrak Q$ is an
isomorphism. Thus $\Gamma_S\to S$ is an isomorphism, proving the universal
mapping property.
\end{proof}

\begin{corollary}[Cohomology of the Poincare bundle]
\label{mumford-cor-poincare-cohomology}
\dcref{M:ch3:III.13}
\uses{mumford-thm-poincare-universal-mapping}
For an abelian variety $X$ of dimension $g$ and its dual $\widehat X$,
\[
H^i(\widehat X\times X,P)=
\begin{cases}0&i\ne g,\\ k&i=g.\end{cases}
\]
\source{mumford-abelian-varieties:page-0140}
\source{mumford-abelian-varieties:page-0141}
\end{corollary}
\begin{proof}
With $S=\operatorname{Spec}k$ and the trivial bundle in the universal
mapping proof, the resulting complex is a free resolution of $k$ over the
regular local ring $\mathcal O_{0,\widehat X}$. Comparing it with the
Koszul resolution on a basis $x_1,\ldots,x_g$ of the maximal ideal modulo
its square gives the asserted cohomology. Explicitly, the Koszul differential
is
\[
d_p(e_{i_1}\wedge\cdots\wedge e_{i_p})=
\sum_{j=1}^p(-1)^j x_{i_j}e_{i_1}\wedge\cdots\widehat e_{i_j}\cdots\wedge e_{i_p}.
\]
\end{proof}

\begin{corollary}[Structure-sheaf cohomology]
\label{mumford-cor-structure-sheaf-cohomology}
\dcref{M:ch3:III.13}
\uses{mumford-cor-poincare-cohomology}
For $\dim X=g$,
\[
\dim_kH^p(X,\mathcal O_X)=\binom gp.
\]
\source{mumford-abelian-varieties:page-0141}
\end{corollary}
\begin{proof}
The cohomology is that of the Koszul complex tensored with $k$; all its
differentials vanish, and the rank of its $p$-cochains is $\binom gp$.
\end{proof}

\begin{corollary}[Tangent space of the dual]
\label{mumford-cor-tangent-dual-h1}
\dcref{M:ch3:III.13}
\uses{mumford-thm-poincare-universal-mapping}
There is a canonical isomorphism
$T_0\widehat X\simeq H^1(X,\mathcal O_X)$.
\source{mumford-abelian-varieties:page-0141}
\source{mumford-abelian-varieties:page-0142}
\end{corollary}
\begin{proof}
For $S=\operatorname{Spec}k[\epsilon]/(\epsilon^2)$, the universal property
identifies $T_0\widehat X$ with the kernel of restriction from line bundles
on $S\times X$ trivial on the closed fibre. The exact sequence
\[
1\to1+\epsilon\mathcal O_X\to\mathcal O_{S\times X}^*\to
\mathcal O_X^*\to1,
\qquad1+\epsilon\mathcal O_X\simeq\mathcal O_X,
\]
gives the result.
\end{proof}

\begin{corollary}[Dual isogenies]
\label{mumford-cor-dual-isogeny-degree}
\dcref{M:ch3:III.13}
\uses{mumford-thm-poincare-universal-mapping,mumford-thm-euler-associated-fibration}
If $f:X\to Y$ is an isogeny, there is a unique homomorphism
$\widehat f:\widehat Y\to\widehat X$ with
\[
(1\times\widehat f)^*P_X\simeq(f\times1)^*P_Y,
\]
and $\deg f=\deg\widehat f$.
\source{mumford-abelian-varieties:page-0141}
\source{mumford-abelian-varieties:page-0142}
\end{corollary}
\begin{proof}
The universal property gives $\widehat f$. Applying the Euler
characteristic formula for a finite flat associated fibre space to the two
pullbacks gives
\[
\chi((1\times\widehat f)^*P_X)=\deg\widehat f\,\chi(P_X),\qquad
\chi((f\times1)^*P_Y)=\deg f\,\chi(P_Y).
\]
Since both Poincare bundles have Euler characteristic $(-1)^g$, the degrees
are equal.
\end{proof}

\begin{corollary}[The morphism associated to a line bundle]
\label{mumford-cor-line-bundle-morphism}
\label{mumford-prop-phi-L}
\dcref{M:ch3:III.13}
\uses{mumford-thm-poincare-universal-mapping,mumford-prop-maximal-triviality-scheme}
For every line bundle $L$ on $X$, $\phi_L:X\to\operatorname{Pic}^0X\simeq
\widehat X$ is a morphism and its scheme-theoretic kernel is $K(L)$.
\source{mumford-abelian-varieties:page-0142}
\end{corollary}
\begin{proof}
The universal property applied to
$m^*L\otimes p_1^*L^{-1}\otimes p_2^*L^{-1}$ gives the morphism. The
definition of $K(L)$ as the maximal subscheme on which this bundle is
trivial is exactly the scheme-theoretic kernel.
\end{proof}

\begin{proposition}[Divisorial correspondences]
\label{mumford-prop-divisorial-correspondence}
\dcref{M:ch3:III.13}
\uses{mumford-thm-poincare-universal-mapping,mumford-thm-euler-associated-fibration}
Let $Q$ be a line bundle on $X\times Y$, with trivial restrictions to the
two coordinate factors, and $\dim X=\dim Y$. The following are equivalent:
\begin{enumerate}
\item no nonzero subscheme of $X$ makes $Q$ trivial;
\item no nonzero subscheme of $Y$ makes $Q$ trivial;
\item $|\chi(Q)|=1$.
\end{enumerate}
When they hold, $Y\simeq\widehat X$ and $X\simeq\widehat Y$.
\source{mumford-abelian-varieties:page-0142}
\source{mumford-abelian-varieties:page-0143}
\end{proposition}
\begin{proof}
The universal property gives $f:Y\to\widehat X$ with
$Q=(1_X\times f)^*P$. Condition (2) says that $f$ has trivial kernel. Equal
dimensions then make $f$ an isogeny, and
$|\chi(Q)|=\deg f\,|\chi(P)|=\deg f$. Conversely, if $f$ has positive
dimensional kernel, finite subgroups of arbitrarily large order $d$ in the
kernel make $Q$ descend through $X\times Y/F$, so the Euler characteristic
is divisible by every such $d$ and is zero. This proves the equivalence.
\end{proof}

\begin{corollary}[Duality hypothesis]
\label{mumford-thm-duality-hypothesis}
\label{mumford-prop-biduality-any-characteristic}
\dcref{M:ch3:III.13}
The canonical morphism $X\to\widehat{\widehat X}$ defined by the Poincare
bundle is an isomorphism.
\source{mumford-abelian-varieties:page-0143}
\uses{mumford-prop-divisorial-correspondence}
\end{corollary}
\begin{proof}
The Poincare bundle satisfies condition (2) of the preceding proposition,
and hence also (1), so the induced morphism is an isomorphism.
\end{proof}

\section{Duality of finite commutative group schemes}
\label{mumford-III14-source}
\dcref{M:ch3:III.14}

\begin{definition}[Cartier dual]
\label{mumford-def-cartier-dual}
\dcref{M:ch3:III.14}
For a finite commutative group scheme $G=\operatorname{Spec}R$ in
characteristic $p>0$, the dual group scheme is
$\widehat G=\operatorname{Spec}R^*$, where the dual of the group-law
coproduct is the multiplication on $R^*$ and the dual of multiplication on
$R$ is its coproduct.
\source{mumford-abelian-varieties:page-0143}
\source{mumford-abelian-varieties:page-0144}
\end{definition}

\begin{theorem}[Cartier dual represents characters]
\label{mumford-thm-cartier-duality}
\dcref{M:ch3:III.14}
\uses{mumford-def-cartier-dual}
For every scheme $S$,
\[
\widehat G(S)\simeq\operatorname{Hom}_S(G_S,\mathbb G_{m,S}),
\]
functorially in $S$.
\source{mumford-abelian-varieties:page-0144}
\source{mumford-abelian-varieties:page-0145}
\source{mumford-abelian-varieties:page-0146}
\end{theorem}
\begin{proof}
For $S=\operatorname{Spec}A$, a homomorphism to $\mathbb G_m$ is an
$\alpha\in R\otimes A$ satisfying
$\mu_R(\alpha)=\alpha\otimes\alpha$ and $\epsilon(\alpha)=1$.
The latter condition implies invertibility, using the inverse map. Since
$R\otimes A\simeq\operatorname{Hom}_A((R\otimes A)^*,A)$, these elements
are precisely algebra homomorphisms $A^*\otimes A\to A$, i.e. $A$-points of
$\operatorname{Spec}R^*$. The construction is natural and respects group
laws. Taking $S=\widehat G$ gives the universal bilinear character
$\widehat G\times G\to\mathbb G_m$.
\end{proof}

\begin{proposition}[Local/reduced decomposition]
\label{mumford-prop-cartier-local-reduced}
\dcref{M:ch3:III.14}
\uses{mumford-thm-cartier-duality}
Every finite commutative group scheme has a unique product decomposition
$G_{r,r}\times G_{r,l}\times G_{l,r}\times G_{l,l}$ according as it and
its Cartier dual are reduced or local. Groups of type $(r,r)$ are products
of cyclic groups of order prime to $p$; groups of type $(r,l)$ are products
of cyclic $p$-groups and type $(l,r)$ are their duals.
\source{mumford-abelian-varieties:page-0147}
\end{proposition}
\begin{proof}
If $G$ is reduced, it is uniquely expressible as a product $G_1\times G_2$,
where $G_1$ has order prime to $p$ and $G_2$ is a $p$-group. By duality,
$\widehat G_1$ is reduced and $\widehat G_2$ is local. When $G$ is local,
splitting $\widehat G$ into local and reduced parts and dualizing back proves
the analogous assertion. The only groups of type $(r,r)$ are reduced groups
of order prime to $p$, hence products of cyclic prime-power groups with
primes distinct from $p$. The only groups of type $(r,l)$ are $p$-groups,
hence products of $\mathbb Z/p^{n_i}\mathbb Z$'s; the groups of type $(l,r)$
are their duals. Applying these decompositions to the four possible types
gives the asserted unique product decomposition.
\end{proof}

\begin{proposition}[Primitive elements and the Lie algebra]
\label{mumford-prop-cartier-primitives}
\dcref{M:ch3:III.14}
\uses{mumford-thm-cartier-duality,mumford-prop-group-scheme-lie}
\[
\operatorname{Hom}_{\mathrm{Gp.Sch}}(\widehat G,\mathbb G_a)
\simeq\{\alpha\in R^*:s^*\alpha=\alpha\otimes1+1\otimes\alpha\}
\simeq\Lie(G).
\]
The $p$-power operation corresponds to composition with Frobenius.
\source{mumford-abelian-varieties:page-0148}
\source{mumford-abelian-varieties:page-0149}
\end{proposition}
\begin{proof}
For $\alpha(1)=0$ and $\alpha(\mathfrak M^2)=0$, the transpose of $\alpha$
is a derivation. Primitive elements are exactly these invariant derivations,
and the coproduct calculation gives the displayed bijections.
\end{proof}

\begin{definition}[Height one]
\label{mumford-def-height-one-group}
\dcref{M:ch3:III.14}
A finite group scheme is of height one if it has one point $e$ and
$x^p=0$ for every $x$ in the maximal ideal at $e$.
\source{mumford-abelian-varieties:page-0150}
\end{definition}

\begin{theorem}[Height-one group schemes and $p$-Lie algebras]
\label{mumford-thm-height-one-equivalence}
\dcref{M:ch3:III.14}
\uses{mumford-prop-cartier-primitives}
The functor $G\mapsto\Lie G$ is an equivalence between finite commutative
height-one group schemes and finite-dimensional $p$-Lie algebras.
\source{mumford-abelian-varieties:page-0150}
\source{mumford-abelian-varieties:page-0151}
\source{mumford-abelian-varieties:page-0152}
\end{theorem}
\begin{proof}
For a $p$-Lie algebra $\mathfrak g$, set
$U(\mathfrak g)=S(\mathfrak g)/\langle\alpha^{(p)}-\alpha^p\rangle$ and
$G(\mathfrak g)=\operatorname{Spec}U(\mathfrak g)^*$. The coproduct
$s(\alpha)=1\otimes\alpha+\alpha\otimes1$ makes this a finite group
scheme of height one. Primitive elements identify its Lie algebra with
$\mathfrak g$. Conversely, the inclusion $\Lie G\to R^*$ extends to a
bialgebra map $U(\Lie G)\to R^*$; its transpose is a group morphism.
The differential is an isomorphism, hence the morphism is a closed immersion
by Nakayama, and both coordinate rings have dimension
$p^{\dim\Lie G}$, so it is an isomorphism.
\end{proof}

\begin{corollary}[Actions of height-one groups]
\label{mumford-cor-height-one-actions}
\dcref{M:ch3:III.14}
\uses{mumford-thm-height-one-equivalence}
For a commutative height-one $G$ with Lie algebra $\mathfrak g$, an action on
$X$ is equivalent to derivations $D_\alpha$ satisfying
\[
[D_\alpha,D_\beta]=0,\qquad D_{\alpha^{(p)}}=D_\alpha^p.
\]
Moreover multiplication by $p$ on $G$ is zero.
\source{mumford-abelian-varieties:page-0152}
\source{mumford-abelian-varieties:page-0153}
\end{corollary}
\begin{proof}
The relations define a homomorphism from $U(\mathfrak g)$ to differential
operators. Transposing it gives
$\mathcal O_X\to R\otimes\mathcal O_X$; the unit and composition diagrams
are exactly the action diagrams. The assertion about multiplication by $p$
follows because it kills the Lie algebra and the height-one equivalence is
faithful.
\end{proof}

\begin{corollary}[Semisimple/nilpotent decomposition]
\label{mumford-cor-p-lie-semisimple-nilpotent}
\dcref{M:ch3:III.14}
\uses{mumford-thm-height-one-equivalence}
Any vector space with a $p$-linear operation decomposes uniquely as
$V=V_s\oplus V_n$, with a basis of fixed vectors on $V_s$ and nilpotent
$p$-operation on $V_n$. For a height-one group this is the decomposition
$\Lie G_{l,r}\oplus\Lie G_{l,l}$.
\source{mumford-abelian-varieties:page-0153}
\source{mumford-abelian-varieties:page-0154}
\end{corollary}
\begin{proof}
For a height-one group scheme, the local/reduced decomposition gives
\[
G=G_{l,r}\times G_{l,l}.
\]
Both factors are again of height one. Since $G_{l,r}$ is killed by $p$, it
is isomorphic to $\mu_p^n$ for some $n$. As $\Lie(\mu_p)$ has a generator
$e$ with $e^{(p)}=e$, the induced decomposition
\[
\Lie G=\Lie G_{l,r}\oplus\Lie G_{l,l}=\mathfrak g_1\oplus\mathfrak g_2
\]
gives a basis $e_1,\ldots,e_n$ of $\mathfrak g_1$ with
$e_i^{(p)}=e_i$. Since $G_{l,l}$ is local and $\mathfrak g_2$ is contained
in the maximal ideal of $\Gamma(\mathcal O_{G_{l,l}})$, every
$\alpha\in\mathfrak g_2$ satisfies $\alpha^{(p^m)}=0$ for $m$ sufficiently
large. Thus the displayed decomposition is invariant under the $p$-linear
operation, with the first summand having a basis of fixed vectors and the
second having nilpotent $p$-operation. The height-one equivalence identifies
these summands with $\Lie G_{l,r}$ and $\Lie G_{l,l}$, respectively; the
uniqueness follows from the corresponding uniqueness of the group-scheme
decomposition.
\end{proof}

\section{Applications to abelian varieties}
\label{mumford-III15-source}
\dcref{M:ch3:III.15}

\begin{theorem}[Kernel of the dual isogeny]
\label{mumford-thm-dual-isogeny-kernel}
\dcref{M:ch3:III.15}
\uses{mumford-thm-cartier-duality,mumford-prop-translation-quotient-isogeny}
If $f:X\to Y$ is an isogeny with kernel $K$ and dual map
$\widehat f:\widehat Y\to\widehat X$ with kernel $K'$, then
$K'\simeq\widehat K$ canonically.
\source{mumford-abelian-varieties:page-0154}
\source{mumford-abelian-varieties:page-0155}
\source{mumford-abelian-varieties:page-0156}
\end{theorem}
\begin{proof}
For a line bundle $L$ on $Y$,
\[
\phi_{f^*L}(x)=\widehat f(\phi_L(f(x))).
\]
Thus $\widehat K(S)$ is the kernel of pullback between families of line
bundles on $S\times Y$ and $S\times X$. Since $S\times Y$ is the quotient
of $S\times X$ by the free action of $K$, these are liftings of the action
to $S\times X\times\mathbb A^1$. Completeness of $X$ makes a linear lift
independent of $x$, so it has the form
\[
k:(s,x,\alpha)\longmapsto(s,x+k,\chi(k,s)\alpha)
\]
for a character $\chi:K\times S\to\mathbb G_m$. Cartier duality identifies
these with $\widehat K(S)$, naturally in $S$.
\end{proof}

\begin{definition}[Height-one isogeny]
\label{mumford-def-height-one-isogeny}
\dcref{M:ch3:III.15}
\uses{mumford-def-height-one-group}
An isogeny $f:X\to Y$ is of height one if $k(X)^p\subset k(Y)$.
This is equivalent to its kernel being a height-one group scheme.
\source{mumford-abelian-varieties:page-0156}
\end{definition}
\begin{proof}
For a height-one isogeny, integral closure gives
$\mathcal O_{Y,0}\supset\{a^p:a\in\mathcal O_{X,0}\}$, so the kernel has
height one. Conversely, invariants of the height-one kernel on affine opens
are exactly the functions coming from $Y$.
\end{proof}

\begin{theorem}[Height-one isogenies and $p$-Lie subalgebras]
\label{mumford-thm-height-one-isogeny}
\dcref{M:ch3:III.15}
\uses{mumford-thm-height-one-equivalence}
Height-one isogenies out of $X$, up to isomorphism, correspond one-to-one to
$p$-Lie subalgebras of $\Lie X$.
\source{mumford-abelian-varieties:page-0157}
\end{theorem}
\begin{proof}
Such isogenies are determined by their height-one kernels, which are
subgroup schemes of $X_p$. The height-one equivalence identifies these with
$p$-Lie subalgebras of $\Lie X_p=\Lie X$.
\end{proof}

\begin{theorem}[Frobenius and the Hasse--Witt operation]
\label{mumford-thm-lie-dual-frobenius}
\dcref{M:ch3:III.15}
\uses{mumford-cor-tangent-dual-h1}
Under $\Lie\widehat X\simeq H^1(X,\mathcal O_X)$, the $p$-power operation
on $\Lie\widehat X$ is the Frobenius map on $H^1(X,\mathcal O_X)$.
\source{mumford-abelian-varieties:page-0158}
\source{mumford-abelian-varieties:page-0159}
\source{mumford-abelian-varieties:page-0160}
\source{mumford-abelian-varieties:page-0161}
\end{theorem}
\begin{proof}
For $D$ represented by an infinitesimal automorphism over
$k[\epsilon]/(\epsilon^2)$, take $p$ copies over
$k[\epsilon_1,\ldots,\epsilon_p]/(\epsilon_i^2)$. Their composite is
$1+s_1D+s_2D^2+\cdots+s_pD^p$. Specializing the elementary symmetric
functions by $s_1\mapsto0$, $s_p\mapsto\epsilon$ gives $D^{(p)}$.
For a group scheme this is translation by the $p$-power of the tangent
point. A tangent vector of the dual represented by a cocycle
$\{1+\epsilon f_{ij}\}$ therefore has $p$-power represented by
$\{1+\epsilon f_{ij}^p\}$, exactly Frobenius on $H^1(X,\mathcal O_X)$.
\end{proof}

\section{Cohomology of line bundles}
\label{mumford-III16-source}
\dcref{M:ch3:III.16}

\begin{theorem}[Riemann--Roch]
\label{mumford-thm-riemann-roch}
\dcref{M:ch3:III.16}
\uses{mumford-prop-degree-hilbert,mumford-thm-cube-varieties}
For $L\simeq\mathcal O_X(D)$ on a $g$-dimensional abelian variety,
\[
\chi(L)=\frac{(D^g)}{g!},\qquad \chi(L)^2=\deg\phi_L.
\]
\source{mumford-abelian-varieties:page-0161}
\source{mumford-abelian-varieties:page-0162}
\source{mumford-abelian-varieties:page-0163}
\end{theorem}
\begin{proof}
Twisting by $\operatorname{Pic}^0X$ does not change Euler characteristic.
For symmetric $L$, $n_X^*(L^m)\simeq L^{mn^2}$, so homogeneity gives
$\chi(L^n)=a(L)n^g/g!$. For very ample $L_i=\mathcal O(D_i)$, the polynomial
$g!\chi(L_1^{n_1}\otimes L_2^{n_2})$ equals
$((n_1D_1+n_2D_2)^g)$; taking $n_1=1,n_2=-1$ yields the first formula.

If $K(L)$ is finite, pullback of the Poincare bundle gives
\[
\chi(m^*L\otimes p_1^*L^{-1}\otimes p_2^*L^{-1})=(-1)^g\chi(L)^2.
\]
The quotient $X\times X\to X\times\widehat X$ has degree $\deg\phi_L$,
so the same expression is $(-1)^g\deg\phi_L$. If $K(L)$ is not finite,
 finite subgroups of arbitrarily large order force both quantities to vanish.
\end{proof}

\begin{remark}[Analytic degree and Pfaffian]
\label{mumford-source-pfaffian-degree}
\dcref{M:ch3:III.16}
Over $\mathbf C$, for $L=L(H,\alpha)$ with $K(L)$ finite,
\[
 \deg\phi_L=\operatorname{order}K(L)
 =\operatorname{order}(U^\perp/U)=\det(E(u_i,u_j)),
 \qquad |\chi(L)|=\sqrt{\det(E(u_i,u_j))}.
\]
Thus $|\chi(L)|$ is the Pfaffian of the skew matrix $(E(u_i,u_j))$ up to
sign, since $\operatorname{pf}(A)^2=\det A$.
\source{mumford-abelian-varieties:page-0165}
\end{remark}

\begin{theorem}[Vanishing and index]
\label{mumford-thm-line-bundle-index}
\dcref{M:ch3:III.16}
\uses{mumford-thm-riemann-roch,mumford-thm-cube-varieties}
If $K(L)$ is finite, there is a unique $i(L)$ with
$H^p(X,L)=0$ for $p\ne i(L)$ and $H^{i(L)}(X,L)\ne0$, and
$i(L^{-1})=g-i(L)$.
\source{mumford-abelian-varieties:page-0161}
\source{mumford-abelian-varieties:page-0164}
\end{theorem}
\begin{proof}
The Cartesian diagram for $1\times\phi_L$ and flatness gives
\[
\dim H^q(X\times X,m^*L\otimes p_1^*L^{-1})=
\begin{cases}0&q\ne g,\\\deg\phi_L&q=g.\end{cases}
\]
Kunneth therefore gives
\[
\sum_{i=0}^qh^i(L)h^{q-i}(L^{-1})=
\begin{cases}0&q\ne g,\\\deg\phi_L&q=g.\end{cases}
\]
Nonnegativity forces exactly one nonzero group on each side and their
indices sum to $g$.
\end{proof}

\begin{lemma}[Polynomial growth of cohomology]
\label{mumford-lem-index-growth}
\dcref{M:ch3:III.16}
For a projective $r$-fold and line bundles $L_1,\ldots,L_k$, there is $c$
such that
\[
\dim H^i(L_1^{a_1}\otimes\cdots\otimes L_k^{a_k})
<c(1+|a|^r).
\]
\source{mumford-abelian-varieties:page-0166}
\source{mumford-abelian-varieties:page-0167}
\source{mumford-abelian-varieties:page-0168}
\end{lemma}
\begin{proof}
Twist by a sufficiently ample bundle to reduce to very ample bundles and
induct on $\dim X+k$. A hyperplane-section exact sequence compares the
cohomology for consecutive exponents with cohomology on a variety of
dimension $r-1$. The inductive bound is $c(1+|a|^{r-1})$; summing over the
exponent gives the asserted $c'(1+|a|^r)$, with negative exponents handled
by the reversed exact sequence.
\end{proof}

\begin{theorem}[Index polynomial]
\label{mumford-thm-index-polynomial}
\dcref{M:ch3:III.16}
If $L$ is ample and $M$ non-degenerate on a $g$-fold, the polynomial
$P(t)=\chi(L^t\otimes M)$ has all roots real, and $i(M)$ is the number of
positive roots counted with multiplicity.
\source{mumford-abelian-varieties:page-0166}
\source{mumford-abelian-varieties:page-0168}
\source{mumford-abelian-varieties:page-0169}
\source{mumford-abelian-varieties:page-0170}
\source{mumford-abelian-varieties:page-0171}
\uses{mumford-lem-index-growth,mumford-thm-line-bundle-index}
\end{theorem}
\begin{proof}
Step A: if $P(1,t)\ne0$ for $0<t<1$, then
$i(L)=i(L\otimes H)$. Otherwise choose the first exponent at which the
index changes. The hyperplane exact sequence gives cohomology on a
dimension $g-1$ section at least $N^gP(1,a/N)$, contradicting the growth
lemma as $N$ tends to infinity.

Step B: for $F(s,t)=\chi(L_1^s\otimes L_2^t)$, nonvanishing of
$F(t,1-t)$ on $(0,1)$ and Step A applied after inserting a very ample $L_3$
give equality of the indices of $L_1$ and $L_2$.

Step C: if all three of $L_1,L_2,L_1\otimes L_2$ are non-degenerate, apply
the map $(x,y)\mapsto x-y$ and Leray to obtain
$i(L_1\otimes L_2)\le i(L_1)+i(L_2)$. Step D compares the indices along a
segment whose Euler polynomial has a unique zero of multiplicity $\lambda$;
the hyperplane exact sequences and the growth lemma give
\(0<i(L_2)-i(L_1)\le\lambda\).

For completeness, the final induction in the source takes hyperplane powers
$H^\alpha$ and proves
\[
H^p(H^\alpha,L_1^k\otimes L_2^l)=0\quad(p<i_2-\alpha-1)
\]
from the exact sequence
\[
0\to L_1^{k-1}\otimes L_2^{l+1}|_{H^{\alpha-1}}
\to L_1^k\otimes L_2^l|_{H^{\alpha-1}}
\to L_1^k\otimes L_2^l|_{H^\alpha}\to0.
\]
The resulting injection and the growth estimate yield
\[
\left|f\left(\frac{r}{N},\frac{s}{N}\right)\right|<\frac{c}{N^{i_2-i_1}},
\qquad
\left|f\left(\frac{r}{N},\frac{s}{N}\right)\right|>\frac{c'}{N^\lambda},
\]
so $i_2-i_1\le\lambda$ and the theorem follows.
\end{proof}

\begin{remark}[Completion of the index root count]
\label{mumford-rem-index-root-count-conclusion}
\dcref{M:ch3:III.16}
\uses{mumford-thm-index-polynomial}
For $P(n,m)=\chi(L^n\otimes M^m)$, choose $q>0$ and integers $r_i$ with
$q^{-1}(r_i-1)<t_i<q^{-1}r_i$, where $t_i$ are the real roots of
$P(t,1)$. By the preceding index comparisons,
\[
i(r_k)\le i(r_k-1)=i(r_{k-1})\le\cdots\le i(r_1-1)=i(0),
\]
and $i(r_i-1)-i(r_i)\le\lambda_i$. For large positive $r$, $i(r)=0$,
and for large negative $r$, $i(r)=g$. Hence
\[
g=i(\text{very negative }r)-i(\text{very positive }r)\le\sum_i\lambda_i.
\]
Since $P$ has degree at most $g$, equality holds; $i(0)=i(M)$ is the sum
of multiplicities of positive roots.
\source{mumford-abelian-varieties:page-0172}
\source{mumford-abelian-varieties:page-0173}
\end{remark}

\begin{corollary}[Complex index and eigenvalues]
\label{mumford-cor-complex-index-negative-eigenvalues}
\dcref{M:ch3:III.16}
\uses{mumford-thm-index-polynomial}
If $X=V/U$ and $L=L(H,\alpha)$ is non-degenerate, then $i(L)$ is the
number of negative eigenvalues of $H$. Choosing an ample
$L_0=L(H_0,\alpha_0)$,
\[
\chi(L_0^n\otimes L)=\operatorname{pf}(nE_0+E),
\qquad E=\operatorname{Im}H.
\]
Diagonalizing relative to $H_0$, $H(e_i,e_j)=\lambda_i\delta_{ij}$, the
Pfaffian is $\prod_i(n+\lambda_i)$, proving the assertion.
\source{mumford-abelian-varieties:page-0173}
\source{mumford-abelian-varieties:page-0174}
\end{corollary}

\section{Very ample line bundles}
\label{mumford-III17-source}
\dcref{M:ch3:III.17}

\begin{theorem}[Very ampleness of powers]
\label{mumford-thm-high-power-very-ample}
\dcref{M:ch3:III.17}
\uses{mumford-thm-index-polynomial,mumford-cor-complex-index-negative-eigenvalues}
For any ample line bundle $L$ on an abelian variety $X$, $L^n$ is very ample
if $n>3$.
\source{mumford-abelian-varieties:page-0174}
\source{mumford-abelian-varieties:page-0175}
\source{mumford-abelian-varieties:page-0176}
\end{theorem}
\begin{proof}
It suffices to prove the assertion for $L^3$. Choose an effective divisor
$D$ without multiple components. For $x,y\in X$,
\[
T_x^*D+T_y^*D+T_{-x-y}^*D\in|3D|.
\]
We establish: (1) for $x_0\ne x_1$, some $D'\in|3D|$ contains $x_0$ but
not $x_1$; (2) every tangent vector at $x_0$ is separated by some
$D'\in|3D|$. By translation assume $x_0=0$.

If (1) failed, then $\operatorname{Supp}D=\operatorname{Supp}D-x_1$, so
$x_1\in K(L)$. Quotienting by the finite group generated by $x_1$ and
comparing dimensions of sections gives a contradiction. Thus (1) holds.

If (2) failed for a nonzero tangent vector $t$, the invariant vector field
$T$ would be tangent to $D$ everywhere. Locally, $T(\phi)=\alpha\phi$.
In characteristic $0$ this forces $t=0$. In characteristic $p$, let
$H\subset K(L)$ be the smallest subgroup containing $t$; $D$ is invariant
under translations by $H$.

Let $X'=X/H$, $D'=D/H$. The quotient $\pi:X\to X'$ is flat and
surjective, $D'$ is a divisor, and $D=\pi^*D'$. Thus
\[
\mathcal O_X(-D)\simeq
\mathcal O_{X'}(-D')\otimes_{\mathcal O_{X'}}\mathcal O_X.
\]
Comparing dimensions of sections on $X$ and $X'$ gives a contradiction.
Therefore (2) holds.
\end{proof}
